% GENERATED FILE. Edit canonical lessons, then run npm run generate:books.
% canonical-source-hash: fnv1a64:cc614de9db5fe918
% canonical-lessons: HI-C255-bahu, HI-C255-damad, HI-C255-devar, HI-C255-nanad, HI-C255-jija

\chapter{A daughter-in-law, A son-in-law, A husband's younger brother, A husband's sister, A sister's husband}
\label{ch:hi-a-daughter-in-law-a-son-in-law-a-husband-s-younger-brother-a-husband-s-sister-a-sister-s-husband}
\hlchaptermodality{255}

\begin{chapteropening}
\textbf{By the end of this chapter:} \emph{I can say a daughter-in-law, a son-in-law, a husband's younger brother, a husband's sister and a sister's husband.}

Complete the last lesson of chapter 255: i can say a daughter-in-law, a son-in-law, a husband's younger brother, a husband's sister and a sister's husband.
\end{chapteropening}

\section[\texorpdfstring{\hi{बहू} (bahū)}{बहू (bahū)}]{\hi{बहू} (bahū) — a daughter-in-law}

\label{lesson:HI-C255-bahu}



\begin{quote}
Before the new one: say the Hindi for a mother-in-law, then the Hindi for a father-in-law.
\end{quote}

\subsection*{You'll want to know: \hi{बहू}}
\textbf{\hi{बहू}} — \emph{bahū} — \textquotedblleft{}a daughter-in-law\textquotedblright{}.

\begin{grammarlens}[title={What you've built}]
One more word for this chapter.
\end{grammarlens}

\subsection*{Your turn}
\begin{itemize}
\raggedright
  \item \emph{Say it:} \emph{bahū}
  \item \emph{Say it:} \emph{bahū}, once more
  \item \emph{Say it:} say \emph{bahū}
  \item \emph{Recall:} say the Hindi for a mother-in-law, then the Hindi for a father-in-law, then say \emph{bahū} again
\end{itemize}

\begin{tcolorbox}[breakable,colback=teal!4,colframe=teal!35!black,title={Before you move on}]
What does \hi{बहू} mean? (\textquotedblleft{}A daughter-in-law\textquotedblright{}.) Say it once more.
\end{tcolorbox}
\section[\texorpdfstring{\hi{दामाद} (dāmād)}{दामाद (dāmād)}]{\hi{दामाद} (dāmād) — a son-in-law}

\label{lesson:HI-C255-damad}



\begin{quote}
Before the new one: say the Hindi for a father-in-law, then the Hindi for a daughter-in-law.
\end{quote}

\subsection*{You'll want to know: \hi{दामाद}}
\textbf{\hi{दामाद}} — \emph{dāmād} — \textquotedblleft{}a son-in-law\textquotedblright{}.

\begin{grammarlens}[title={What you've built}]
One more word for this chapter.
\end{grammarlens}

\subsection*{Your turn}
\begin{itemize}
\raggedright
  \item \emph{Say it:} \emph{dāmād}
  \item \emph{Say it:} \emph{dāmād}, once more
  \item \emph{Say it:} say \emph{dāmād}
  \item \emph{Recall:} say the Hindi for a father-in-law, then the Hindi for a daughter-in-law, then say \emph{dāmād} again
\end{itemize}

\begin{tcolorbox}[breakable,colback=teal!4,colframe=teal!35!black,title={Before you move on}]
What does \hi{दामाद} mean? (\textquotedblleft{}A son-in-law\textquotedblright{}.) Say it once more.
\end{tcolorbox}
\section[\texorpdfstring{\hi{देवर} (devar)}{देवर (devar)}]{\hi{देवर} (devar) — a husband's younger brother}

\label{lesson:HI-C255-devar}



\begin{quote}
Before the new one: say the Hindi for a daughter-in-law, then the Hindi for a son-in-law.
\end{quote}

\subsection*{You'll want to know: \hi{देवर}}
\textbf{\hi{देवर}} — \emph{devar} — \textquotedblleft{}a husband's younger brother\textquotedblright{}.

\begin{grammarlens}[title={What you've built}]
One more word for this chapter.
\end{grammarlens}

\subsection*{Your turn}
\begin{itemize}
\raggedright
  \item \emph{Say it:} \emph{devar}
  \item \emph{Say it:} \emph{devar}, once more
  \item \emph{Say it:} say \emph{devar}
  \item \emph{Recall:} say the Hindi for a daughter-in-law, then the Hindi for a son-in-law, then say \emph{devar} again
\end{itemize}

\begin{tcolorbox}[breakable,colback=teal!4,colframe=teal!35!black,title={Before you move on}]
What does \hi{देवर} mean? (\textquotedblleft{}A husband's younger brother\textquotedblright{}.) Say it once more.
\end{tcolorbox}
\section[\texorpdfstring{\hi{ननद} (nanad)}{ननद (nanad)}]{\hi{ननद} (nanad) — a husband's sister}

\label{lesson:HI-C255-nanad}



\begin{quote}
Before the new one: say the Hindi for a son-in-law, then the Hindi for a husband's younger brother.
\end{quote}

\subsection*{You'll want to know: \hi{ननद}}
\textbf{\hi{ननद}} — \emph{nanad} — \textquotedblleft{}a husband's sister\textquotedblright{}.

\begin{grammarlens}[title={What you've built}]
One more word for this chapter.
\end{grammarlens}

\subsection*{Your turn}
\begin{itemize}
\raggedright
  \item \emph{Say it:} \emph{nanad}
  \item \emph{Say it:} \emph{nanad}, once more
  \item \emph{Say it:} say \emph{nanad}
  \item \emph{Recall:} say the Hindi for a son-in-law, then the Hindi for a husband's younger brother, then say \emph{nanad} again
\end{itemize}

\begin{tcolorbox}[breakable,colback=teal!4,colframe=teal!35!black,title={Before you move on}]
What does \hi{ननद} mean? (\textquotedblleft{}A husband's sister\textquotedblright{}.) Say it once more.
\end{tcolorbox}
\section[\texorpdfstring{\hi{जीजा} (jījā)}{जीजा (jījā)}]{\hi{जीजा} (jījā) — a sister's husband}

\label{lesson:HI-C255-jija}



\begin{quote}
Before the new one: say the Hindi for a husband's younger brother, then the Hindi for a husband's sister.
\end{quote}

\subsection*{You'll want to know: \hi{जीजा}}
\textbf{\hi{जीजा}} — \emph{jījā} — \textquotedblleft{}a sister's husband\textquotedblright{}.

\begin{grammarlens}[title={What you've built}]
One more word for this chapter.
\end{grammarlens}

\subsection*{Your turn}
\begin{itemize}
\raggedright
  \item \emph{Say it:} \emph{jījā}
  \item \emph{Say it:} \emph{jījā}, once more
  \item \emph{Say it:} say \emph{jījā}
  \item \emph{Recall:} say the Hindi for a husband's younger brother, then the Hindi for a husband's sister, then say \emph{jījā} again
\end{itemize}

\begin{tcolorbox}[breakable,colback=teal!4,colframe=teal!35!black,title={Before you move on}]
What does \hi{जीजा} mean? (\textquotedblleft{}A sister's husband\textquotedblright{}.) Say it once more.
\end{tcolorbox}
